\subsection{SQL Injection 与 Remote Code Execution：传统漏洞被 LLM 放大}

slides 31、35--37 把传统 SQL injection 和 RCE 放入 Agentic Hybrid System。传统漏洞来自未转义输入进入解释器；Agent 场景增加了一个生成层，攻击文本可能先诱导 LLM 产生恶意 query 或 code，再由工具执行。即使攻击者无法直接调用数据库，只要能影响模型上下文并让模型拥有高权限工具，仍可能完成间接利用。

\lecturefigure{slides-images/slide-031.png}{官方 slide 31：传统 SQL Injection 与 Agent attack chain 的起点}
\lecturefigure{slides-images/slide-035.png}{官方 slide 35：Agentic Hybrid System 中 SQL Injection 的完整链路}
\lecturefigure{slides-images/slide-036.png}{官方 slide 36：Remote Code Execution 的传统与 Agent 路径}
\lecturefigure{slides-images/slide-037.png}{官方 slide 37：RCE attack chain 的完整状态}

\begin{importantbox}{参数化与权限隔离仍然有效}
LLM 时代没有取消传统 secure coding。数据库使用 parameterized query，shell 尽量不用字符串拼接，工具参数使用 schema validation，执行账户采用最小权限，写操作放入事务并提供补偿。模型 hardening 可以减少恶意输出概率，却不能替代这些确定性边界。
\end{importantbox}

\begin{warningbox}{不要把“模型拒答率”当 RCE 防线}
只要攻击成功概率非零，高频调用就会累积风险；攻击者还会换语言、编码、上下文和间接载体。系统应假设危险字符串终会出现，并保证即使出现，也无法越过 capability boundary。
\end{warningbox}

